\section{3.1 Why Whole-Repository Analysis Is Required}
The final garage pipeline cannot be explained only from \texttt{garage\_lab\_combined}. Critical assumptions, scripts, and design choices emerged across multiple stages. This chapter maps those stages and clarifies folder-level responsibilities.

\subsection{Table 3.1 Whole-repository folder-to-function map}
\begin{verbatim}
| Folder | Role in project evolution | Key artifacts |
|---|---|---|
| `src/calibration`, `src/core`, `src/legacy` | Early algorithm prototypes and baseline 3D logic | `triangulate_3d.py`, `render_3d_robot.py`, `main_3d_tracker.py` |
| `src/capture`, `src/tools`, `scripts` | Utility capture and calibration automation | `auto_sport_calibrate.py`, `record_dual_cameras.py` |
| `GARAGE_CAMERAS` | Practical multi-camera recording reliability layer | `record_cams.py`, `sync_record_2.py` |
| `garage-20260217T113109Z-3-001` | Imported baseline arena calibration and environment scripts | `extrinsics_1`, `environment/reconstruction.py` |
| `garage_lab_combined` | Unified final pipeline and GT protocols | `process_4cam_to_3d.py`, `gt_eval/*`, `thesis/*` |
| `output`, `data`, `cal`, `Intrinsicsdec17` | Historical data, calibration snapshots, and rendered evidence | videos, calibration JSONs |
\end{verbatim}

\section{3.2 Stage A: Baseline Geometry and Rendering (\texttt{src/*})}
The earliest phase built foundational capabilities: video capture, 2D detection integration, geometric triangulation, and 3D rendering. \texttt{src/core/triangulate\_3d.py} and \texttt{src/core/triangulate\_v2.py} implemented multi-view reconstruction logic with projection-based consistency checks. \texttt{src/core/render\_3d\_robot.py} and related scripts provided visual sanity checks for skeleton and ball trajectories.

At this stage, the system could demonstrate concept feasibility but lacked robust operational procedures. Scripts were often tuned to specific sessions and had limited standardized evaluation.

\section{3.3 Stage B: Capture Robustness Under Real Constraints (\texttt{GARAGE\_CAMERAS})}
As practical experiments intensified, capture became a bottleneck. USB camera IDs changed after reconnects, frame rates varied, and operator workflow required preview plus recording control. The \texttt{GARAGE\_CAMERAS} folder addressed these constraints with ffmpeg/openCV recording tools, camera discovery helpers, and synchronized recording scripts.

This stage established an important lesson: poor capture discipline can invalidate even strong triangulation algorithms. Reliable input streams are a prerequisite for reproducible geometry.

\section{3.4 Stage C: External Baseline Transfer (\texttt{garage-20260217...})}
This folder supplied a prior garage setup with established AprilTag geometry and reconstruction scripts. It included:
\begin{itemize}
  \item \texttt{extrinsics\_1}: camera pose and arena visualization workflow,
  \item \texttt{environment}: dual-camera high-speed tracking utilities,
  \item \texttt{Intrinsics}: board/intrinsics references.
\end{itemize}

The transfer was not direct reuse. Intrinsics/extrinsics from the old setup were useful for understanding but not always valid for the current camera placement and lens conditions. The major outcome was methodological: reuse structure, recalibrate parameters.

\section{3.5 Stage D: Unified Pipeline Construction (\texttt{garage\_lab\_combined})}
\texttt{garage\_lab\_combined} became the operational research package integrating:
\begin{itemize}
  \item config management (\texttt{config/cameras.yaml}),
  \item intrinsics/extrinsics scripts,
  \item synchronized short-clip recorder,
  \item 4-camera processing to 3D JSON,
  \item rendering and live views,
  \item GT evaluation pipelines for ball and joints.
\end{itemize}

This stage transformed the project from loosely coupled scripts into a coherent scientific workflow.

\section{3.6 Stage E: Quantitative Validation and Bias Modeling (\texttt{gt\_eval})}
The \texttt{gt\_eval} subtree formalized evaluation protocol design and reporting. For ball localization, static and dynamic protocols were separated to distinguish geometric bias from temporal instability. For joints, touch-based protocols highlighted the difference between algorithmic precision and human placement uncertainty.

This stage enabled rigorous statements such as "raw mean error is 150.77 mm" instead of subjective statements such as "trajectory looks better".

\section{3.7 Cross-Folder Dependency Graph}
A simplified dependency graph is:

\begin{enumerate}
  \item \texttt{GARAGE\_CAMERAS} and \texttt{garage\_lab\_combined/scripts/record\_*} produce synchronized clips.
  \item \texttt{garage\_lab\_combined/cal/*} produces intrinsics and extrinsics JSON files.
  \item \texttt{process\_4cam\_to\_3d.py} consumes clips + calibration and outputs 3D motion JSON.
  \item render scripts consume motion JSON + arena map and generate videos/plots.
  \item evaluation scripts consume motion JSON + trials CSV and output reports/correction models.
\end{enumerate}

The historical \texttt{src/*} and \texttt{garage-20260217.../*} layers provide the mathematical and practical foundations for this final graph.

\section{3.8 Repository-Level Contribution}
This cross-folder evolution is a contribution in itself. In many student projects, older scripts are abandoned and final results become hard to reproduce. Here, the lineage is preserved and can be audited: from first prototype to quantitative evaluation and thesis-grade documentation.


\section{3.9 Lessons Learned from Legacy-to-Unified Transition}
The transition from legacy scripts to unified workflow produced several engineering lessons:

\begin{enumerate}
  \item Naming conventions matter.
\end{enumerate}
   Early scripts used generic camera names (\texttt{Cam1}, \texttt{Cam2}) while the unified system uses physical-role names (\texttt{camNorth}, \texttt{camEast}, \texttt{camSouth}, \texttt{camWest}). This reduced class of mistakes where processing used correct files with wrong physical semantics.

\begin{enumerate}
  \item Sessionized outputs are mandatory.
\end{enumerate}
   Moving from ad-hoc output filenames to timestamped session folders made experiment repeatability and report traceability significantly better.

\begin{enumerate}
  \item Visual diagnostics should be first-class outputs.
\end{enumerate}
   Overlay and 3D render diagnostics were not decorative; they detected calibration failure earlier than scalar metrics alone.

\begin{enumerate}
  \item Partial recalibration is operationally valuable.
\end{enumerate}
   In realistic operation, one camera may move while others remain stable. Being able to recalibrate and merge one camera's extrinsics reduced downtime.

\section{3.10 Data Governance and Reproducibility Practices}
Repository maturity improved when the team adopted data governance habits:
\begin{itemize}
  \item preserve raw clips before trimming/alignment,
  \item keep processing outputs immutable per session,
  \item store trial CSV and metadata near result files,
  \item never overwrite reference calibration files without versioning,
  \item generate summary reports as machine-readable JSON and human-readable markdown.
\end{itemize}

This thesis benefits directly from these practices because every major claim can be traced to a concrete file.

\section{3.11 How Historical Folders Continue to Add Value}
Although \texttt{garage\_lab\_combined} is the active integration folder, historical folders remain useful:
\begin{itemize}
  \item \texttt{src/*} preserves algorithmic prototypes and alternative renderers,
  \item \texttt{garage-20260217...} preserves independent calibration/reporting approaches,
  \item \texttt{GARAGE\_CAMERAS} keeps robust standalone capture tools.
\end{itemize}

Maintaining these layers avoids knowledge loss and supports troubleshooting when regressions appear in the main pipeline.



\section{3.12 Detailed Chronology of Technical Milestones}
The following chronology clarifies how engineering decisions were driven by observed failure modes rather than by theoretical preference:

\begin{itemize}
  \item Initial semester phase: early stereo and 3D rendering scripts demonstrated that reconstructed trajectories were possible from low-cost cameras, but produced unstable scale and orientation between sessions.
  \item Lab phase: more controlled calibration and multi-camera processing produced convincing demonstration outputs (\texttt{final\_3d\_robot.mp4}), but assumptions from that environment did not transfer directly to the garage.
  \item Garage migration phase: adoption of arena tags and room dimensions enabled physical-world interpretation, revealing hidden inconsistencies that were previously masked in lab-only visualizations.
  \item Calibration hardening phase: repeated intrinsics and extrinsics recalibration, plus overlay-based diagnostics, improved geometric reliability.
  \item Protocol phase: formal GT campaigns replaced impression-based evaluation and identified exact bias structure in ball and joint estimates.
\end{itemize}

This chronology is important because it shows methodological maturity: from prototype confidence to evidence-driven iteration.

\section{3.13 Data Artifacts and Their Role in Scientific Argumentation}
For each claim in this thesis, supporting artifacts exist:
\begin{itemize}
  \item Calibration quality claims: intrinsics JSON, extrinsics JSON, overlay image sets.
  \item Tracking claims: per-session motion JSON files and rendered videos.
  \item Accuracy claims: summary metrics and per-trial CSV reports from GT scripts.
  \item Robustness claims: dynamic metrics from slow/fast/no-ball scenarios.
\end{itemize}

The artifact-oriented workflow allowed rapid dispute resolution when contradictory observations appeared. For example, when a render looked physically implausible, it was possible to inspect camera maps, overlays, and reprojection metrics to isolate root cause.

\section{3.14 Integrating Human Workflow into Technical Design}
A recurring practical issue was operator distance from workstation during capture. This influenced script evolution toward:
\begin{itemize}
  \item automatic triggers based on corner quality,
  \item settle delays before recording,
  \item compact keyboard controls and status overlays,
  \item deterministic output naming.
\end{itemize}

This is not trivial UI work; it directly affects data quality and therefore scientific validity. In research systems involving physical spaces, user workflow is a technical dependency.

\section{3.15 Repository Sustainability for Future Students}
The project now has a structure that future researchers can extend:
\begin{itemize}
  \item clear folder responsibilities,
  \item documented session pipelines,
  \item reproducible report generators,
  \item thesis-linked artifact paths.
\end{itemize}

This sustainability is a meaningful contribution in academic engineering labs where toolchains often become unusable after one cohort.
